\section*{Problem 233 --- Round 2}

\subsection*{1) ROUND-2 OBJECTIVE}
Path \textbf{(C): obstruction/correction}.  The conjectured bound
\[
\sum_{n\le N} d_n^2 \ll N(\log N)^2
\]
remains open.  Round~1 had only the lower bound and a trivial upper bound through the maximal gap.  In this round I do two strictly new things:
\begin{itemize}
\item I prove the exact equivalence between the $N$-formulation and the standard $x$-formulation
\[
\sum_{p_n\le x} d_n^2 \ll x\log x;
\]
\item I import the current unconditional mean-square theorem of Stadlmann and translate it carefully into an explicit bound in the original $N$-variable.
\end{itemize}

\subsection*{2) Round-1 FOUNDATION USED}
I use the Round~1 notation $d_n=p_{n+1}-p_n$ and the Round~1 lower bound
\[
\sum_{n\le N} d_n^2\ge \frac{(p_{N+1}-2)^2}{N}\gg N(\log N)^2,
\]
which I do not reprove here.

\subsection*{3) NEW INSIGHT / TOOL (ROUND-2)}
The new ingredients are:
\begin{itemize}
\item the exact identity
\[
M(x):=\sum_{p_n\le x} d_n^2 = \sum_{n\le \pi(x)} d_n^2;
\]
\item Stadlmann's theorem that
\[
\frac{1}{\pi(x)}\sum_{p_n\le x} d_n^2 = O_{\varepsilon}(x^{23/100+\varepsilon})
\qquad (\varepsilon>0),
\]
which implies
\[
M(x)\ll_{\varepsilon} \frac{x^{123/100+\varepsilon}}{\log x}.
\]
\end{itemize}

\subsection*{4) ATTACK PLAN (ROUND-2)}
The exact Round~1 gap was the absence of any nontrivial upper bound in the original variable $N$.  So I first prove that the conjecture is equivalent to the $x$-formulation.  Then I plug the current best unconditional $x$-bound into this equivalence to obtain a genuine unconditional theorem about
\[
S(N):=\sum_{n\le N} d_n^2.
\]

\subsection*{5) WORK (ROUND-2)}
Define
\[
M(x):=\sum_{p_n\le x} d_n^2.
\]
If $\pi(x)=N$, then the condition $p_n\le x$ is equivalent to $n\le N$, so
\[
M(x)=\sum_{n\le N} d_n^2.
\]
In particular,
\[
M(p_N)=\sum_{n\le N} d_n^2 = S(N).
\]
This makes the two formulations equivalent.

\medskip
\noindent\textbf{Equivalence lemma.}
The following are equivalent:
\begin{itemize}
\item[(i)] $S(N)\ll N(\log N)^2$;
\item[(ii)] $M(x)\ll x\log x$.
\end{itemize}

\noindent\emph{Proof.}
Assume (i).  Let $N=\pi(x)$.  Then $M(x)=S(N)\ll N(\log N)^2$.  By the prime number theorem,
\[
N\asymp \frac{x}{\log x}
\qquad\mbox{and}\qquad
\log N\asymp \log x,
\]
so $M(x)\ll x\log x$.  Thus (ii) follows.

Conversely, assume (ii).  Then
\[
S(N)=M(p_N)\ll p_N\log p_N.
\]
Again by the prime number theorem, $p_N\asymp N\log N$, hence
\[
S(N)\ll N(\log N)^2.
\]
So (i) follows.  This proves the equivalence.

\medskip
Now apply Stadlmann's theorem:
\[
\frac{M(x)}{\pi(x)}=O_{\varepsilon}(x^{23/100+\varepsilon}).
\]
Multiplying by $\pi(x)$ and using $\pi(x)\ll x/\log x$, we get
\[
M(x)\ll_{\varepsilon} \frac{x^{123/100+\varepsilon}}{\log x}.
\]
Set $x=p_N$.  Then
\[
S(N)=M(p_N)\ll_{\varepsilon} \frac{p_N^{123/100+\varepsilon}}{\log p_N}.
\]
Using $p_N\ll N\log N$ and $\log p_N\asymp \log N$, we conclude that
\[
S(N)\ll_{\varepsilon} N^{123/100+\varepsilon}(\log N)^{23/100+\varepsilon}.
\]
Therefore Round~2 proves the following corrected theorem.

\medskip
\noindent\textbf{Corrected theorem proved in Round~2.}
For every $\varepsilon>0$,
\[
\sum_{n\le N} d_n^2 \ll_{\varepsilon} N^{123/100+\varepsilon}(\log N)^{23/100+\varepsilon}.
\]
Combining this with the Round~1 lower bound gives the rigorous sandwich
\[
N(\log N)^2 \ll \sum_{n\le N} d_n^2 \ll_{\varepsilon} N^{123/100+\varepsilon}(\log N)^{23/100+\varepsilon}.
\]

\subsection*{6) ADVERSARIAL VERIFICATION}
\begin{itemize}
\item \emph{Any indexing error in $M(p_N)$?} No.  Since $p_n\le p_N$ iff $n\le N$, we have exactly
\[
M(p_N)=\sum_{n\le N} d_n^2.
\]
\item \emph{Did the argument use only valid asymptotic input?} Yes: only the prime number theorem bounds $\pi(x)\asymp x/\log x$ and $p_N\asymp N\log N$.
\item \emph{Could the Stadlmann theorem have used a different normalization?} The cited source states explicitly that the average considered is
\[
\frac{1}{\pi(x)}\sum_{p_n\le x}(p_{n+1}-p_n)^2,
\]
which is exactly the normalization used here.
\item \emph{Does this settle the conjecture?} No.  The exponent of $N$ is still $1.23+\varepsilon$ rather than $1$, so the main barrier remains substantial.
\end{itemize}

\subsection*{7) FINAL}
\textbf{UNRESOLVED (BUT STRICTLY ADVANCED).}  The conjectured bound $\sum_{n\le N} d_n^2\ll N(\log N)^2$ is still open, but Round~2 replaces the Round~1 trivial upper bound by the current best unconditional type of estimate in the original variable:
\[
\sum_{n\le N} d_n^2 \ll_{\varepsilon} N^{123/100+\varepsilon}(\log N)^{23/100+\varepsilon}.
\]
It also proves the exact equivalence between the $N$-formulation and the standard $x$-formulation.

\subsection*{8) COMPLETION ESTIMATE}
COMPLETION: 35\%

\subsection*{9) REFERENCES}
\begin{itemize}
\item J. Stadlmann, \emph{On the mean square gap between primes}, arXiv:2212.10867.
\item J. Stadlmann, \emph{Topics in sieve theory}, D.Phil. thesis, University of Oxford, 2024.
\end{itemize}
